%! Multiplying \& Dividing Rational Expressions — Unit 4, Lesson 4.2 — Exit Ticket (Answer Key)
\documentclass[10pt]{article}
\usepackage{algebra2-article}
\usepackage{algebra2-key}   % pulls in -boxes; do NOT also load -boxes
\newcommand{\TallMath}[1]{$\displaystyle #1\rule[-1.4em]{0pt}{3.2em}$}

\begin{document}

\pageheader{Unit 4, Lesson 4.2}{Exit Ticket --- Answer Key}
\namedateperiod

\vspace{0.08in}

No notes. Show the factored line for question 1.

\vspace{0.06in}

\begin{enumerate}[leftmargin=1.8em, itemsep=9pt]
  \item \textbf{Simplify and restrict.} \; $\dfrac{x^2-9}{x^2+2x}\div\dfrac{x-3}{x+2}$

        \vspace{4pt}
        After keep--change--flip, factored: \ans{$\frac{(x-3)(x+3)}{x(x+2)}\cdot\frac{x+2}{x-3}$}

        \vspace{4pt}
        Simplest form: \ans{$\frac{x+3}{x}$} \quad All restrictions: $x\neq$ \ans{$0,\ -2,\ 3$}

        \vspace{4pt}
        Two of your restrictions are invisible in the simplest form. Name them and say where each
        came from. \\[3pt]
        \ans{$x=-2$} came from \ans{the original denominator $x^2+2x$} \\[3pt]
        \ans{$x=3$} came from \ans{the divisor's numerator ($x-3=0$)}

  \item \textbf{Explain.} A classmate says: ``Once you flip, $x+2$ is on top, so $x=-2$ isn't a
        restriction anymore.'' Is the classmate right? Explain what flipping does and does not change.
        \ansline{No. The restrictions belong to the problem \emph{as written}, and there $x+2$ is a denominator:}
        \ansline{at $x=-2$ the very first fraction has no value, so there is nothing to flip in the first place.}
        \ansline{Flipping changes how the expression \emph{looks}, not which inputs it accepts --- and it hides a restriction rather than removing one.}

  \item \textbf{SOL-style multiple choice.} Which is the correct simplest form of
        $\dfrac{x^2-25}{x^2+3x}\div\dfrac{x-5}{x+3}$, with \emph{all} of its restrictions?
        \begin{enumerate}[label=\Alph*., leftmargin=1.9em, itemsep=1pt]
          \item $\dfrac{x+5}{x}$, \; $x\neq0$
          \item $\dfrac{x+5}{x}$, \; $x\neq0,\,-3$
          \item $\dfrac{x+5}{x}$, \; $x\neq0,\,-3,\,5$
                \quad \textcolor{keyred}{\textbf{$\leftarrow$ correct}}
          \item $\dfrac{x}{x+5}$, \; $x\neq0,\,-3,\,5$
        \end{enumerate}
        \vspace{2pt}
        {\small $\dfrac{(x-5)(x+5)}{x(x+3)}\cdot\dfrac{x+3}{x-5}=\dfrac{x+5}{x}$, with $x\neq0,-3$ from
        the two denominators and $x\neq5$ from the divisor's numerator. \textbf{A} reads the
        restrictions off the \emph{answer}, \textbf{B} forgets the divisor's numerator (today's
        signature error), and \textbf{D} flipped the \emph{first} fraction instead of the divisor.}
\end{enumerate}

\begin{teachernote}
Graded for completion --- ``mistakes happen, blanks don't.'' Sort into four piles: \textbf{(1) got it};
\textbf{(2) missing source 3} --- chose \textbf{B}, or listed only $x\neq0,-2$ in item 1; \textbf{(3)
restrictions from the answer} --- chose \textbf{A}, which is Lesson 4.1's error still unresolved;
\textbf{(4) flipped the wrong fraction} --- chose \textbf{D}, a mechanical slip that a single worked
example fixes. Pile 2 is the one that matters: the divisor's numerator is the only genuinely new idea
today, and it returns in 4.3 (a complex fraction is a division, so its bottom fraction's numerator is
restricted) and again in 4.6, where a solution that lands on \emph{any} of these values is extraneous.
Item~2 earns full credit only with both halves: the restriction is read from the problem as written,
\emph{and} flipping hides it rather than removing it.
\end{teachernote}

\end{document}
